% =========================================================
% The Logarithm — Notes
% Chapter: elementary-functions
% Section: logarithmic
% Volume III · Analysis
% =========================================================

\section{The Logarithm}
\label{sec:logarithmic}

% ---------------------------------------------------------
% Toolkit
% ---------------------------------------------------------

\begin{tcolorbox}[title={Toolkit --- Logarithmic Functions}]
\begin{itemize}
  \item $\ln : (0,\infty) \to \mathbb{R}$: inverse of $e^x$.
    Equivalently, $\ln x = \int_1^x 1/t\, dt$.
  \item Functional equation: $\ln(xy) = \ln x + \ln y$.
    Converts products to sums.
  \item $(\ln x)' = 1/x$;\; $\ln 1 = 0$;\; $\ln e = 1$.
  \item Change of base: $\log_a x = \ln x / \ln a$.
  \item Fundamental logarithmic limit:
    $\lim_{x \to 0}\ln(1+x)/x = 1$.
  \item Growth dominance: for all $r > 0$,
    $(\ln x)/x^r \to 0$ as $x \to \infty$ and
    $x^r \ln x \to 0$ as $x \to 0^+$.
\end{itemize}
\end{tcolorbox}

% ---------------------------------------------------------
% Definition — Natural Logarithm
% ---------------------------------------------------------

\begin{tcolorbox}[title={Definition --- Natural Logarithm}]
\begin{definition}[Natural Logarithm]
\label{def:natural-logarithm}
The \textbf{natural logarithm} $\ln : (0,\infty) \to \mathbb{R}$ is
the inverse of the exponential function: for $x > 0$,
\[
  \ln x \;:=\; \text{the unique } y \in \mathbb{R}
  \text{ such that } e^y = x.
\]
Equivalently, $\ln x$ is defined by the integral
\[
  \ln x \;:=\; \int_1^x \frac{1}{t}\, dt.
\]
\end{definition}
\end{tcolorbox}

\begin{remark*}[Standard quantified statement]
\[
  \forall x > 0 :\; e^{\ln x} = x.
  \qquad
  \forall y \in \mathbb{R} :\; \ln(e^y) = y.
\]
\end{remark*}

\begin{remark*}[Two definitions, same function]
The inverse-of-exponential definition requires $e^x$ to be
established first. The integral definition is self-contained and
is preferred when building analysis from scratch. Both yield the
same function; the choice depends on what has been proved at the
time of introduction.
\end{remark*}

\begin{remark*}[Key properties]
\begin{enumerate}[label=(\roman*)]
  \item $\ln 1 = 0$;\; $\ln e = 1$.
  \item \textbf{Functional equation:} $\ln(xy) = \ln x + \ln y$
    for all $x, y > 0$.
  \item $\ln(x/y) = \ln x - \ln y$;\;
    $\ln(x^r) = r\ln x$ for $r \in \mathbb{R}$.
  \item $\ln$ is strictly increasing on $(0,\infty)$.
  \item $\ln x \to +\infty$ as $x \to \infty$;\;
    $\ln x \to -\infty$ as $x \to 0^+$.
  \item $(\ln x)' = 1/x$ for all $x > 0$.
\end{enumerate}
\end{remark*}

\begin{remark*}[General logarithm]
For $a > 0$, $a \neq 1$, the \textbf{logarithm base $a$} is
$\log_a x := \ln x / \ln a$. Then $\log_a(a^x) = x$ and
$a^{\log_a x} = x$. The common logarithm is $\log_{10}$; the
binary logarithm is $\log_2$.
\end{remark*}

% ---------------------------------------------------------
% Proposition — Fundamental Logarithmic Limit
% ---------------------------------------------------------

\begin{proposition}[Fundamental Logarithmic Limit]
\label{prop:limit-logarithm-fundamental}
\[
  \lim_{x \to 0} \frac{\ln(1+x)}{x} = 1.
\]

\smallskip
\noindent
\hyperref[prf:limit-logarithm-fundamental]{\textit{Go to proof.}}
\end{proposition}

\begin{remark*}[Standard quantified statement]
\[
  \forall\varepsilon > 0,\;
  \exists\delta > 0 :\;
  0 < |x| < \delta \Rightarrow
  \left|\frac{\ln(1+x)}{x} - 1\right| < \varepsilon.
\]
\end{remark*}

\begin{remark*}[Equivalence with exponential limit]
The fundamental exponential limit $\lim(e^x-1)/x = 1$ and the
fundamental logarithmic limit $\lim\ln(1+x)/x = 1$ are equivalent:
each follows from the other via the substitution $x \leftrightarrow
e^x - 1$ (which maps $x \to 0$ to $x \to 0$). Both express the
first-order Taylor approximation of the respective function at $0$.
\end{remark*}

\begin{remark*}[Proof strategy]
From the series: $\ln(1+x) = x - x^2/2 + x^3/3 - \cdots$ (valid for
$|x| < 1$). Dividing by $x$ gives $1 - x/2 + x^2/3 - \cdots \to 1$
as $x \to 0$ by continuity of power series.
\end{remark*}

% ---------------------------------------------------------
% Proposition — Every Power Dominates the Logarithm
% ---------------------------------------------------------

\begin{proposition}[Every Power Dominates the Logarithm]
\label{prop:power-dominates-log}
For every $r > 0$:
\[
  \lim_{x \to \infty} \frac{\ln x}{x^r} = 0
  \qquad\text{and}\qquad
  \lim_{x \to 0^+} x^r \ln x = 0.
\]

\smallskip
\noindent
\hyperref[prf:power-dominates-log]{\textit{Go to proof.}}
\end{proposition}

\begin{remark*}[Standard quantified statement]
\[
  \forall r > 0 :\;
  \lim_{x \to \infty} x^{-r} \ln x = 0
  \;\wedge\;
  \lim_{x \to 0^+} x^r \ln x = 0.
\]
\end{remark*}

\begin{remark*}[Proof strategy]
For the first: substitute $x = e^t$ to get $t/e^{rt} \to 0$ by
exponential dominance. For the second: substitute $x = 1/t$ and
reduce to the first limit: $x^r \ln x = (-\ln t)/t^r \to 0$.
\end{remark*}

\begin{remark*}[Significance]
These limits establish the position of $\ln x$ in the growth
hierarchy: $\ln x$ is dominated by every positive power. Combined
with the exponential dominance result, the full ordering is
$\ln x \ll x^r \ll e^x$ as $x \to \infty$ for any $r > 0$.
The second limit says that near $0^+$, the positive power $x^r$
wins over the logarithmic blow-up $|\ln x| \to \infty$. This is
used in integrability arguments near $0$.
\end{remark*}
